\documentclass[11pt,a4paper]{article}
\usepackage[T1]{fontenc}
\usepackage{lmodern,amsmath,amssymb,microtype,booktabs}
\usepackage[margin=25mm]{geometry}
\usepackage[colorlinks=true,urlcolor=blue,linkcolor=blue]{hyperref}
\newcommand{\ga}{\mathfrak a}\newcommand{\gb}{\mathfrak b}
\newcommand{\Rh}{\widehat R}
\title{Compact-transfer extension and the physical\\two-creation limit in dS$_3$}
\author{Alan Day \and Chengheng\thanks{AI research assistant (OpenAI).}}
\date{Working research note --- 31 October 2026}
\begin{document}\maketitle
\begin{abstract}
A cosh product removes the small-numerator restriction in our preceding
crossing estimate. Cutting off the actual channel vectors, and bounding
their omitted Hilbert norms uniformly in the observer mass, then removes
the compact cutoff after the mass limit. Within the specified regulated
$G_N=0$ model this yields the different-species crossed limit and the
two-creation Bose free-observer baseline. We retain a failed numerical
reference, reproduce its near-pole precision loss, and provide both
exact-label and increased-precision controls. These calculations neither
establish finite-gravity dynamics nor complete the full Fock construction.
\end{abstract}

\section{Model and the statement to be justified}
Keep $\ell=1$, $G_N=0$, proper $SO^+(3,1)$, two distinguishable BD
scalar observers with principal-series labels at least a fixed $s>0$,
KG matter labels $p,q>0$, and unit local coupling.
All creation blocks use the previous half-period regulation.
Static time is generated by $h(x)=\sqrt{1+x^2}$.
Condition the initial label to $c=M$, by a limit of normalized narrow
mass packets. Put
\[
 b=M+u,\qquad a=M+u+v,\qquad d=M+v+r.
\]
The conditional physical measure is $da\,db\,dd=du\,dv\,dr$.
No additional Plancherel factor is inserted. Energy coordinates would
instead supply the Jacobian $\prod_{x=a,b,d}h(x)/x$.

Let $U=U_{00}(b,d;a,q,c,p)$ have the preceding physical CG-period phase.
It is a scalar matrix element of the full recoupling map, not an
independently unitary scalar transform. On every fixed compact set of
$(u,v,r,p,q)$ with positive $p,q$, we show
\begin{equation}
 U=\frac{\sin(r\log M-\delta(r))}{\pi r}+O_{\mathcal K}(M^{-1}),
 \label{kernel}
\end{equation}
with the continuous value at zero and the previous four-Gamma phase
$\delta$. The resulting compact distributional limit acts on $dr$.
We then prove a mass-uniform Hilbert tail bound and take $M\to\infty$
before removing that cutoff. The model result for identical matter is
\[
 N_{B,1}(t),\,N_{B,2}(t),\,\langle B_1(t),B_2(t)\rangle
 \longrightarrow L_p(0)^2+|L_p(t)|^2,\qquad D_B(t)\longrightarrow0,
\]
uniformly on compact real times; $L_p$ is given in Section 5.
The proof uses the established model normalization \cite{prior}
and known CG/Racah theory \cite{ds,lprs}; it is a self-checked working
derivation, with no claim of priority.

\newpage
\section{An arbitrary-compact coefficient bound}
Write $A_j=1/2+i\alpha_j$, with
\[
 \alpha=\tfrac12(v+q,p-u,-u-p,v-q),\quad z=-ir/2,\quad |r|\le R.
\]
Set $A_*=\max|\alpha_j|+R/2$. There is now no condition $A_*\le1/2$.
The cosh product \cite{products} gives, for every real $\alpha$,
\begin{equation}
 \frac{|(1/2+i\alpha)_n|}{n!}
 =\frac{(1/2)_n}{n!}
 \sqrt{\prod_{k=0}^{n-1}\left(1+\frac{\alpha^2}{(k+1/2)^2}\right)}
 \le\sqrt{\frac{\cosh(\pi\alpha)}{n+1}}.
 \label{poch}
\end{equation}
Indeed $((1/2)_n/n!)^2\le1/(n+1)$ follows by induction since
$(n+1)^3-(n+1/2)^2(n+2)=(3n+2)/4>0$ for $n\ge0$.
The infinite product bounds the finite one by positivity.

In the LPRS/Wilson decomposition, $X=c_AF+c_BG$ and
$Y=y_AF+y_BH$, where
\[
 F={}_4F_3(A-z;1+id,1-ib-z,1-z;1),\quad
 H={}_4F_3(A;1+i(M+v)-z,1-ib,1+z;1).
\]
The notation $A-z$ shifts all four numerators.
Let $B=\max(|u|,|v|)+R$, $m=M-B>0$, and $C_{\rm ser}=\cosh^2(\pi A_*)$.
For their terms, \eqref{poch} implies
\[
 |f_n|,|h_n|\le\frac{C_{\rm ser}}{(n+1)^2}
       \prod_{j=1}^n\frac{j^2}{j^2+m^2},\qquad
 |f_n'|\le6H_n|f_n|,\quad |h_n'|\le H_n|h_n|.
\]
Primes denote $r$ derivatives; $H_n=\sum_{j=1}^n1/j$.
The derivative estimates use only real parts $1/2$ and one, so
enlarging the bounded imaginary shifts does not change them.

Put $\ga=zy_A$, $\gb=-zy_B$, $S=(\ga-\gb)/z$, with derivative
values at zero. Both coefficients equal one at zero and $F(0)=H(0)$.
The regular Gamma formulas from the preceding note give
$|\ga|,|\gb|\le C_0$ and $|\ga'|+|\gb'|\le C_0Q$, where
\begin{align*}
 C_0&=\cosh(\pi A_*)\sqrt{1+R/(2m)}\,
          e^{\pi R/4}(1-e^{-2\pi m})^{-1/2},\\
 P_1(t)&=\gamma+1+\operatorname{arsinh}t,\\
 P_{1/2}&=\gamma+2\log2+2+\operatorname{arsinh}(2A_*),\\
 Q&=P_1(R/2)+2P_1(M+B)+2P_{1/2}.
\end{align*}
These follow from Gamma moduli and
$|\psi(\sigma+it)|\le|\psi(\sigma)|+\sigma^{-1}
+\operatorname{arsinh}(|t|/\sigma)$ for $\sigma>0$.
The constants grow with the chosen compact set; they are not global
bounds uniform in the transfers.

\newpage
\section{Merged error, full companion and the transverse limit}
Merge before dividing: $Y=(\ga F-\gb H)/z$.
For every term the numerator vanishes at zero. Integrating its derivative
and using $\sum_{n\ge1}H_n/(n+1)^2\le2$ gives
\begin{align}
 |F-1|&\le R_F=\frac{C_{\rm ser}(\zeta(2)-1)}{1+m^2},\nonumber\\
 |Y-S|&\le R_Y=
 \frac{2C_0C_{\rm ser}\{(\zeta(2)-1)Q+14\}}{1+m^2}.
 \label{merged}
\end{align}
Termwise differentiation is justified by the displayed summable bound.
The merged series through $n=N-1$ has exact-arithmetic tail
\begin{equation}
 T_Y(N)=\frac{2C_0C_{\rm ser}}{N}
 \prod_{j=1}^N\frac{j^2}{j^2+m^2}\{Q+7(2+\log N)\}.
 \label{tailseries}
\end{equation}

For completeness, $G={}_4F_3(A+ib;1+i(b+M+v+r/2),1+ib,
1+i(b-r/2);1)$. Set $\beta=b-|r|/2$ and
$A_\beta=\max|\alpha_j|+|r|/2$. The previous term-ratio argument
requires $\beta>0$ and $b+M+v+r/2\ge\beta$, both true for large $M$.
It gives
\[
 |g_n|\le B_\beta=\sqrt{\frac{\sinh\pi\beta}{\pi\beta}}\,
       e^{2\pi A_\beta+\pi A_\beta^2/\beta}.
\]
The exact squared-ratio polynomial is monic of degree ten.
If $N_*=2+\lceil\max_{j<10}|p_j|\rceil$ for its lower coefficients,
the polynomial and all its derivatives are positive at $N_*$.
Thus $|G|\le B_G=(3N_*+3)B_\beta$.
Each $p_j$ has degree at most eight in the bounded-shift mass $M$,
so $N_*=O_{\mathcal K}(M^8)$.
No small-$A_*$ condition was used in this argument.
Gamma prefactors give $|E_0|B_G=O_{\mathcal K}(M^8e^{-\pi M/2})$,
where $D_0=C_+c_A$ and $E_0=C_+c_B$.

Set $Y_*=2C_0Q+R_Y$ and
$D_*=C_0(4\pi)^{-1}\sqrt{\sinh(\pi R/2)/(\pi R/2)}$.
The full physical kernel satisfies
\begin{equation}
 |U-2\Re(D_0S)|\le
 2D_*(R_Y+R_FY_*)+2|E_0|B_GY_* .
 \label{full}
\end{equation}
The remainder is $O_{\mathcal K}(\log M/M^2)$ plus the exponential term.

To pass from Gamma-exact $D_0S$ to \eqref{kernel}, use shifted
Stirling and digamma remainders \cite{stirling}.
All small Gamma arguments have real part $1/2$ or one; all large
ones stay in fixed sectors away from the negative real axis.
Consequently the preceding $C^1$ coefficient proof extends to any
fixed compact shifts. More explicitly, after extracting $M^{\pm z}$,
the two coefficient errors $E_A,E_B$ have $C^1$ norm $O_{\mathcal K}(M^{-1})$
and share $E_A(0)=E_B(0)=\epsilon_0=O_{\mathcal K}(M^{-2})$.
The latter is the peak-phase cancellation, an algebraic identity in
the shifts. Hence
\[
 \frac{E_A(r)e^{-ir\log M}-E_B(r)}z
 =O_{\mathcal K}(M^{-1}+\log M/M^2).
\]
This proves the uniform estimate through zero, rather than a
punctured estimate containing $1/|r|$.

\newpage
\section{Removing the physical spectral cutoff}
The phase in \eqref{kernel}, continuous from $\delta(0)=0$, is
\[
 \delta(r)=\Im\left[2\log\Gamma(1-ir/2)+
       \sum_j\{\log\Gamma(A_j+ir/2)-\log\Gamma(A_j)\}\right].
\]
The preceding Dirichlet argument therefore applies on every fixed
compact interval: for uniformly bounded $C_c^2$ tests $f_M$,
$\int f_M U\,dr=f_M(0)+O(1/\log M)+O(1/M)$.
Uniformity also holds over compact $(u,v,p,q)$.
The full physical amplitude is
\[
 A_M=\frac{g(a)}{g(M)}
 \sqrt{\Rh_q(a,b)\Rh_p(b,M)\Rh_p(a,d)\Rh_q(d,M)},\quad
 g(x)=\frac{x e^{\pi[h(x)-x]}}{1+e^{-2\pi x}}.
\]
Here $\Rh_p(x,y)=(A-B)/[8\pi(A+C)(B+C)]$, with
$A=\cosh\pi(x+y)$, $B=\cosh\pi(x-y)$, $C=\cosh\pi p$.
The prior weight estimate gives locally in $C_r^2$, uniformly for
$|t|\le T$,
\[
 A_M e^{it(h_a+h_M-h_b-h_d)}
 =(1+(u+v)/M)F_\infty+O_{\mathcal K,T}(M^{-2}),
\]
\[
 F_\infty=\sqrt{k_q(v)k_p(u)k_p(u-r)k_q(v+r)}\,e^{-itr},
 \qquad k_p(w)=\frac1{8\pi(\cosh\pi w+\cosh\pi p)}.
\]

The actual $s$-channel vector has squared density
$\frac{g(a)}{g(M)}\Rh_q(a,b)\Rh_p(b,M)$ with respect to $da\,db$.
Since $\Rh_p\le k_p$, it is bounded by
\[
 \frac{B_s}{16\pi^2}
 \left(1+\frac{|u|+|v|}{s}\right)e^{-\pi(|u|+|v|)},\qquad
 B_s=2e^{\pi/(\sqrt{1+s^2}+s)},\quad M\ge s.
\]
The other channel has the same bound in its own two transfers.
Choose smooth $0\le\eta_K\le1$, equal to one for $|x|\le K$ and
zero for $|x|\ge2K$. Multiply the channel vectors respectively by
$\eta_K(a-M)\eta_K(b-M)$ and $\eta_K(a-M)\eta_K(d-M)$.
Their product in the crossed integral is precisely
\[
 \chi_K=\eta_K(u+v)^2\eta_K(u)\eta_K(v+r).
\]
It has compact support $|u|\le2K$, $|v|\le4K$, $|r|\le6K$.
Each deleted squared norm is bounded by
\begin{align}
 \epsilon_K&=\frac{B_se^{-\pi K}}{4\pi^2}
 \left[\frac K\pi+\frac1{\pi^2}
 +\frac1s\left(\frac{K^2}\pi+\frac{2K}{\pi^2}+\frac2{\pi^3}\right)\right],
 \label{epsilon}\\
 \|\psi\|^2&\le\overline N_s=\frac{B_s}{4\pi^4}+\frac{B_s}{2s\pi^5}.
 \nonumber
\end{align}
Indeed deletion implies $|u|+|v|>K$; integration over four quadrants
reduces to $\int_K^\infty x(1+x/s)e^{-\pi x}dx$.

\newpage
\section{Full pairings, Bose interference and escaping preparations}
The full CG change of basis is unitary, including its spin channels
and the correct reflection representatives \cite{ds}. Its scalar
compression is a contraction. Cauchy--Schwarz therefore yields
\begin{equation}
 |X_M(t)-X_M^K(t)|\le
 2\sqrt{\epsilon_K\overline N_s},
 \qquad M\ge s,\quad t\in\mathbb R.
 \label{hilbert}
\end{equation}
This concerns integrals of channel vectors. It does not assert the
pointwise inequality $|U|\le1$, or absolute integrability of $|A_MU|$.
The full pairing is defined by the Hilbert-space limit of these cutoffs.

For each fixed $K$, the compact theorem and physical $C^2$ weight
estimate apply. Since $F_\infty(u,v,0)=k_p(u)k_q(v)$,
\[
 X_M^K(t)\longrightarrow
 I_K=\int_{\mathbb R^2}\eta_K(u+v)^2\eta_K(u)\eta_K(v)
                 k_p(u)k_q(v)\,du\,dv .
\]
Dominated convergence gives $I_K\to L_p(0)L_q(0)$, where
\begin{equation}
 L_p(t)=\int_{\mathbb R}k_p(w)e^{-itw}\,dw
       =\frac{\sin(pt)}{4\pi\sinh(\pi p)\sinh t},\qquad
 L_p(0)=\frac{p}{4\pi\sinh(\pi p)}.
 \label{L}
\end{equation}
More explicitly, take the $\limsup_{M\to\infty}$ in
\[
 \sup_{|t|\le T}|X_M(t)-L_p(0)L_q(0)|
 \le2\sqrt{\epsilon_K\overline N_s}
   +\sup_{|t|\le T}|X_M^K(t)-I_K|
   +|I_K-L_p(0)L_q(0)|,
\]
then send $K\to\infty$.
The noncrossed norms already have the same limit.
Thus the different-species ordering difference tends to zero in this
two-creation sector.

For identical matter let $k=T\alpha$, $l=T\beta$ be the previous
labelled vectors, $\alpha_b=e^{-it(h_a-h_b)}$,
$\beta_b=e^{-it(h_b-h_M)}$, and let $\mathsf F$ exchange the matter slots.
The Bose vectors are $\sqrt2P_+k,\sqrt2P_+l$, with
$P_+=(1+\mathsf F)/2$. Define
\[
 N=\|k\|^2=\|l\|^2,\quad Y=\langle k,l\rangle,\quad
 E_1=\langle k,\mathsf Fk\rangle,\quad
 E_2=\langle l,\mathsf Fl\rangle,\quad Z=\langle k,\mathsf Fl\rangle .
\]
The exchange phases are $e^{\pm it(h_d-h_b)}$, which at $r=0$
tend to $e^{\pm it(v-u)}$; the crossed phase of $Z$ tends to one.
The identical cutoff and tail proof, and direct dominated convergence,
give
\[
 N,Z\to L_p(0)^2,\qquad Y,E_1,E_2\to |L_p(t)|^2.
\]
Therefore $N+E_i$ and $Y+Z$ give the stated Bose baseline, and
$D_B=2N+E_1+E_2-2\Re(Y+Z)\to0$. The factors are fixed by the standard
Fock symmetrization, not by an additional factorial.

The result also holds after averaging over any normalized probabilities
on $[s,\infty)$ escaping to infinity: their probability of $[s,C]$
must tend to zero for every fixed $C$. Split the average at $C$, use
the uniform norm bound below $C$ and the conditional limit above it.
No moment assumption is needed. A fixed nonmoving packet does not
satisfy this hypothesis.

\newpage
\section{Numerical checks, retained failures and remaining work}
We evaluated 18 complete kernels at 90 decimal digits:
$M=16,64,256$, $r=-0.6,0,0.6$, $R=1$, and the two families
\[
 (u,v,p,q)=(1.4,-0.8,1.2,1.7),\qquad(-0.9,1.1,2.2,0.6).
\]
Both are outside the former small-numerator window.
All satisfy \eqref{full}, separately allowing their saved truncation
tails. The merged cutoff is $N=160$. These are finite-point tests,
not new numerical quadratures of the full physical observable.

Six unmerged 120-digit references initially gave five passes and one
failure, retained in \texttt{study.json}. At $M=256$, the second family
and $r=0$, the stable result is approximately $1.88534345914$,
while the original reference gives $1.87011200511$:
absolute difference $0.0152314540$, relative difference $0.00807887$.
Its claimed series tail is about $4.79\,10^{-134}$, which does not
include rounding.

The driver parsed transfers at 90 digits, then separately serialized
the computed labels at 120 digits. Exact rational reconstruction of
those strings gives $a+c-b-d=10^{-117}$ instead of zero.
Thus the reference misses the removable-pole branch and subtracts
two huge terms with a recorded condition number about $3.32\,10^{110}$.
The following controls preserve that failure:
\begin{center}
\begin{tabular}{lrr}
\toprule
Input and working precision & collision branch & difference from stable value\\
\midrule
Same rounded strings, 120 digits & no & $1.5231\,10^{-2}$\\
Same rounded strings, 210 digits & no & $4.9324\,10^{-88}$\\
Exact rational labels, 120 digits & yes & $4.9323\,10^{-88}$\\
Exact rational labels, 170 digits & yes & $4.9323\,10^{-88}$\\
\bottomrule
\end{tabular}
\end{center}
The comparison allowance remains $10^{-65}$.
This diagnoses arithmetic cancellation rather than modifying the
stable data or declaring the original reference correct.
None of these controls is an interval rounding certificate.

An independent Decimal/Fraction replay checks the enlarged constants,
full complex assembly, finite error bound, exact companion polynomial
derivatives, input hashes, reference differences and Hilbert tails.
Its initial run passed 441 of 442 checks; the failure was a structural
SymPy comparison of equal factored and distributed polynomials.
A separate refined audit uses the exact zero-polynomial difference
and passes 444 checks without changing data or tolerances.
The 20 finite Pochhammer checks supplement the proof of \eqref{poch};
they do not establish it by sampling.
All failed outputs and prior scientific files remain unchanged.

This result covers regulated two-creation pairings at $G_N=0$,
fixed $s>0$ and compact real times. It does not include the intermediate
vacuum contractions of four self-adjoint fields, a complete Fock
domain, moving thresholds, logarithmically growing times, or finite
gravity. It supplies an iterated-limit argument, not an optimized
global rate in $M$: the compact constants grow with $K$.
The next task is to establish the full operator-domain construction
and its vacuum contractions before comparing finite-$G_N$ dynamics.

\newpage
\section*{Provenance and reproducibility}
The extension uses the known cosh and sinh products, Gamma asymptotics,
CG completeness, and the LPRS crossing formula.
The new work recorded here is their application to the specified
physical channel conventions and cutoff removal.
No exhaustive novelty claim is made.

The implementation is \path{code/ds3_compact_global.py}.
Its driver is \path{code/run_ds3_compact_global.py}.
The dedicated collision diagnosis retains rounded and exact labels,
working precisions, complex values, tails and condition numbers.
The original and refined independent audits are separate files.
Results, dependencies, hashes and compilation review are under
\path{results/stage4/ds3-compact-global/}.
The repository also contains the earlier normalizations and symbolic
phase certificates on which the proof relies.
The standalone PDF is compiled from this saved source and checked
page by page; its input/output hashes accompany the release.

\begin{thebibliography}{9}
\bibitem{prior} Alan Day and Chengheng, project working notes:
\emph{dS$_3$ two-matter sector}, \emph{Physical crossed packets},
\emph{Same-species Bose normalization},
\emph{Complete crossing peak limit},
\emph{Transverse phase}, \emph{Uniform transverse limit}, and
\emph{Physical spectral weights}, 31 October 2026.
The present note extends their explicitly restricted statements;
it does not rewrite the historical files.
\bibitem{ds} S. E. Derkachov and V. P. Spiridonov,
\href{https://arxiv.org/html/1711.07073v2}
{On the $6j$-symbols for $\mathrm{SL}(2,\mathbb C)$ group},
arXiv:1711.07073v2, Section 3.1, orthogonality and completeness,
and Section 4, Racah change of basis.
\bibitem{lprs} J. Liu, E. Perlmutter, V. Rosenhaus and D. Simmons-Duffin,
\href{https://arxiv.org/html/1808.00612v3}
{$d$-dimensional SYK, AdS Loops, and $6j$ Symbols},
arXiv:1808.00612v3, Eqs.\ (3.35)--(3.38).
\bibitem{wilson} W. Groenevelt,
\href{https://arxiv.org/html/math/0306424v1}{The Wilson function transform},
arXiv:math/0306424v1, Propositions 4.4 and 4.7.
Wilson duality and the removable-pole Dirichlet argument are precedents,
not newly introduced mechanisms.
\bibitem{products} NIST DLMF,
\href{https://dlmf.nist.gov/4.36}{Section 4.36},
Eqs.\ 4.36.1 and 4.36.2, sinh and cosh products.
\bibitem{stirling} NIST DLMF,
\href{https://dlmf.nist.gov/5.11}{Section 5.11},
log-Gamma, digamma and shifted Gamma-ratio asymptotics.
Web sources checked 31 October 2026.
\end{thebibliography}
\end{document}
